The capabilities of language models and the systems built around them are developing rapidly. Consequently, the results in this thesis represent one set of models and configurations at a particular point in time, rather than a final assessment of language-model-based robot collaboration. Model choice, decoding parameters, prompt design, context construction, and tool access can all affect the resulting behaviour. Future evaluations should therefore treat these choices as experimental variables and document them sufficiently to distinguish improvements in the underlying model from improvements in the surrounding system.

One direction is to extend the planner with retrieval-augmented generation and external tools exposed through Model Context Protocol servers. Retrieval could provide task-relevant information without placing all available knowledge in every prompt, while reusable skills could define how the model accesses simulator state, robot capabilities, operational rules, or previous experience. Such additions should be evaluated separately from the base model because they alter the information and actions available to the planner. Of particular interest is whether they improve grounded decision-making and generalisation to new warehouse configurations, rather than only increasing system complexity.

The shared-context architecture studied here gives each agent access to a common state description, but it does not investigate communication as a behaviour in its own right. A further study could allow agents to communicate their local state, intended action, or request for assistance directly to one another. With a retained communication history or an explicit adaptation mechanism, repeated interactions could then be examined to determine whether the agents converge towards a useful communication protocol: what information is sent, when communication occurs, and which agents need to receive it. Comparing unrestricted communication with bandwidth-, latency-, or cost-constrained communication would help establish whether the exchanged messages contribute to collaboration or merely duplicate information already available in the environment state.

Deployment-oriented model optimisation is another open direction. Tools such as Embedl\footnote{\url{https://www.embedl.com/}},
which target the optimisation of models for specific hardware platforms,
may reduce latency, memory use, or energy consumption and thereby make local inference more practical. However, techniques used to obtain these gains may also change model behaviour. It would therefore be useful to evaluate the unmodified and optimised versions of a model on the same scenarios, measuring not only execution performance but also action validity, planning quality, robustness, and collaborative task performance. This would reveal whether hardware-specific optimisation introduces a capability loss that is relevant to robot decision-making~\cite{dawane2026onddeviceslm}.

The experiments reported in this thesis were executed on an NVIDIA RTX A4000. Repeating the same protocol on other GPU generations, accelerator architectures, and resource-constrained edge hardware would test the portability of the findings. The comparison should keep the model artifact, inference configuration, and simulator conditions fixed where possible, while recording differences in numerical precision, runtime, memory use, and energy consumption. This would help determine whether the observed performance trends are stable across hardware or partly dependent on the inference stack used in the present study.

Finally, the limited number of completed runs and simulator initialisations restricts the statistical confidence of the present analysis. Future experiments should use more runs, multiple independently selected seeds, and balanced coverage of every model--scenario--configuration combination. Reporting uncertainty intervals and sensitivity across seeds would make it possible to distinguish persistent trends from variation caused by initial conditions or stochastic generation. Useful extensions would also include stricter structured-output validation, comparison with multi-agent reinforcement-learning baselines, and a joint analysis of task performance, inference latency, token usage, and energy cost.
